\chapter*{Advertencia preliminar}
\addcontentsline{toc}{chapter}{Advertencia preliminar}
\markboth{Advertencia preliminar}{Advertencia preliminar}

Este libro continúa \textit{El dispositivo salteño} sobre un territorio de dos municipios. El apéndice E de aquel libro dejó cuatro hipótesis abiertas sobre el mercado inmobiliario como dimensión económica propia del dispositivo --- numeradas H.20 a H.23 --- y admitió que su documentación excedía lo que un autor sin infraestructura de investigación colectiva podía emprender. Este volumen no las confirma. Las lleva a una escala donde son verificables, y en el camino descubre que la hipótesis con la que empezaron era incorrecta.

Aquel apéndice suponía un mercado de suelo que opera contra un Estado ausente. Lo que este relevamiento encontró es otra cosa, y son dos.

La primera: dos regímenes estatales de producción de suelo que funcionan sobre el mismo territorio, con normas propias, organismos propios y poblaciones destinatarias distintas, y cuya diferencia central no es la legalidad sino la intensidad con que el Estado se ocupa de cada uno --- una intensidad inversamente proporcional a la capacidad de sus destinatarios.

La segunda apareció sola, por repetición, en el propio trabajo de conseguir los documentos: la opacidad de este territorio no está en el acto sino en el criterio. Seis instrumentos de dos niveles de gobierno ---el municipal y el provincial--- publican la envoltura y reservan a un anexo el contenido que permitiría verificar si alguien fue tratado con la misma vara que otro; otros ---al menos siete, cinco del régimen de tierra fiscal, uno de Recursos Hídricos y uno de Seguridad--- lo publican sólo como imagen escaneada, que ningún buscador encuentra. \textbf{Seis es lo que quedó de un inventario sistemático de trece ---depurado cuando ocho de sus casos resultaron publicados y ampliado con uno hallado al verificarlos---; no es el total}: cada documento incorporado después de cerrar ese recuento agregó casos de la misma especie, incluido el propio Código de Ordenamiento Territorial del municipio, que declara no urbanizables ambas márgenes del río y remite el ancho de esa franja a un plano que sólo se publica como imagen, sin escala que permita medirlo. Nada de ello es ilegal, y por eso funciona.

A esas dos se sumó una tercera cuando el relevamiento retrocedió hasta el primer número del Boletín Oficial, en 1908: las unidades territoriales, los nombres y las formas institucionales son los mismos desde hace más de un siglo. El mercado inmobiliario no construyó este dispositivo, lo encontró armado --- aunque no intacto, porque el municipio había perdido en el camino la voz que alguna vez tuvo sobre el agua de su jurisdicción.

Ninguna de las tres es una tesis sobre captura. El capítulo~\ref{cap:planteo} las enuncia, los capítulos~\ref{cap:fincas} y~\ref{cap:siglo} documentan la tercera y el~\ref{cap:opacidad} desarrolla la segunda.

Conviene contestar acá, y no en la página seiscientos, la pregunta que este material provoca
antes que ninguna otra: \textbf{si esto es una denuncia de corrupción}. No lo es, y el motivo no
es prudencia sino resultado. En todo lo relevado ---los años 1908 a 1921 del Boletín Oficial leídos hoja
por hoja, mil cuatrocientas noventa y dos ediciones entre 1923 y 1945, los años
1946 a 1950, un expediente judicial, dos códigos de ordenamiento y una década de
resoluciones--- \textbf{este libro no encontró un solo acto ilegal, ni un pago, ni una
instrucción, ni un acto dirigido a beneficiar a alguien con nombre}. Cada uno de los actos que
produce el resultado que el libro describe tiene su fundamento propio y su autoridad
competente.

Tampoco es la historia de un malentendido. Lo que hay es un conjunto de omisiones que, durante
más de un siglo, beneficiaron siempre al mismo lado: criterios que no se publican, líneas que se
publican dos días y no quedan asentadas en la matrícula, un municipio al que se le asignó
regular el suelo sin darle con qué, y una intensidad regulatoria que baja cuando sube la
capacidad del destinatario. \textbf{Ese resultado no necesita que nadie lo haya querido, y eso
es lo que lo vuelve difícil de corregir}: no se arregla encontrando a un culpable. El capítulo
\ref{cap:conclusion} lo dice al final con los documentos ya en la mano, y el
\ref{cap:prospectiva} explica por qué la palabra del título no equivale a conspiración ni a
accidente.

Cuatro advertencias sobre el registro.

\textbf{Primera.} La escala municipal agrava la exigencia metodológica en lugar de aliviarla. En un departamento de unos doce mil habitantes todos se conocen, y el costo de nombrar recae sobre personas identificables. La regla del libro anterior --- afirmaciones sobre posiciones estructurales, nunca sobre individuos morales --- rige aquí con más rigor. Un intendente que aprueba un loteo defectuoso no es, por ese hecho, un caso de corrupción; es, hasta que un expediente diga otra cosa, un caso de capacidad técnica insuficiente frente a una presión de mercado desproporcionada. El capítulo~\ref{cap:hacienda} muestra que esa lectura es, además, la correcta.

\textbf{Segunda.} Las cajas grises que aparecen a lo largo del texto marcan tipográficamente el punto exacto donde la exposición pasa de lo que consta en documento a lo que el autor infiere. Nada inferido sale de esa caja. Cuando una afirmación se apoya en prensa y no en instrumento público, se dice.

\textbf{Tercera.} \textbf{El estándar de cobertura de este libro no es uniforme, y conviene saberlo antes de leer cualquier afirmación sobre lo que falta.} Buena parte del Boletín Oficial se barrió por término de búsqueda sobre el texto extraído de los PDF, el período 1951--2012 sólo se consultó de manera puntual, y los años 1947 a 1950 se barrieron enteros pero sin lectura sobre la imagen (capítulo~\ref{cap:metodo}): ahí, que algo no aparezca significa solamente que el barrido no lo encontró. Hay treinta y dos tramos leídos edición por edición y verificados sobre la imagen ---de los treinta y seis que releva el apéndice \ref{cap:fuentes}, quedan fuera los cuatro de 1947 a 1950---, que forman dos ventanas: las \textbf{196 primeras ediciones del Boletín}, del número 1 del 15 de octubre de 1908 al 196 de octubre de 1910, a las que se suman \textbf{el resto de 1910 y los años 1911, 1912, 1913 y 1914 completos}; y las \textbf{mil quinientas tres} que van del número 962 al 2.464, junio de 1923 a diciembre de 1945 ---de las que quedan once sin leer, seis ausentes y cinco en blanco, que el apéndice \ref{cap:fuentes} enumera una por una: \textbf{mil cuatrocientas noventa y dos leídas}. A ellas se suman, al cierre, \textbf{el año 1946 completo} ---280 ediciones, del número 2.465 al 2.744--- y los años \textbf{1915} y \textbf{1916}, que en el momento de incorporarse eran el tramo más incompleto y el más limpio de la serie: 1915 con 49 ediciones de la 537 a la 586, sin la 581 y sin las de enero y la primera semana de febrero; 1916 con 49 ediciones de la 587 a la 635, sin un solo salto y con las 196 hojas presentes; \textbf{1917}, 47 ediciones de la 636 a la 683 sin la 640, 289 hojas de las 329 del año y sin OCR propio; \textbf{1918}, 52 ediciones de la 684 a la 736 ---falta entera la 703, que el repositorio reemplaza por un aviso de que no la encontró---, 534 hojas de los 587 folios del año, también sin OCR propio; \textbf{1919}, 49 ediciones de la 738 a la 786, 440 hojas de los 478 folios, con 38 hojas ausentes que caen todas al final de veintiocho ediciones y ninguna en el medio de ninguna; \textbf{1920}, 48 ediciones de la 788 a la 839 ---faltan cuatro: la 793, la 822, la 837 y la 838---, 526 hojas distintas y seiscientos folios; \textbf{1921}, 43 ediciones de la 840 a la 888 ---faltan seis---, 742 hojas y la cadena de folios del 1.681 al 2.532, que empalma sin hueco con el cierre de 1920; y \textbf{el primer semestre de 1922}, 27 ediciones de la 890 a la 916 y 446 hojas, del 13 de enero al 21 de julio, \textbf{único tramo que no cubre un año entero} ---de la edición 917 en adelante no hay archivo---, sin OCR propio, con un techo de resolución de 150 a 192 ppi, sin edición el viernes 14 de abril y con unas trece hojas ausentes dentro de las ediciones entregadas, que la cadena de folios detecta y el conteo de archivos no; y \textbf{1923 entero}, 52 ediciones de la 939 a la 990 y 673 hojas, una por cada viernes del año, sin un solo hueco en la serie ---el tramo más completo de la serie hasta que se incorporó 1924---, también sin OCR propio, con un techo de 150 ppi, con quince hojas que la cadena de folios declara ausentes hasta la edición 970 y con 322 hojas sin control de folio desde la 971, porque la numeración reinicia en cada edición; y \textbf{1924 entero}, 52 ediciones de la 991 a la 1042 y 792 hojas, \textbf{el primer año de toda la serie en que el calendario cierra solo} y en el que el cotejo entre el sumario y las páginas citadas cierra en las 52 ediciones; y \textbf{1925 entero}, 52 ediciones de la 1043 a la 1094 y 798 hojas, también sin huecos, con el control de paginación comprobado sobre veinticuatro hojas interiores y no sobre todas; y \textbf{1926 y 1927 enteros}, 52 ediciones cada uno ---1096 a 1147 y 1148 a 1199---, 769 y 745 hojas, las dos series semanales completas y sin huecos, con las cadenas de saldos incompletas y declaradas como tales; y \textbf{1928 entero}, 52 ediciones de la 1200 a la 1251, con dos meses sin resumen de Tesorería y con el salto de numeración de decretos más grande de la serie, del 7097 al 8001 en un solo día; y \textbf{1929, 1930, 1931 y 1932 enteros} ---ediciones 1252 a 1461---, los cuatro sin OCR propio, con un techo de 150 ppi, sin ninguna salida de barrido automático y \textbf{sin un solo plano publicado en los cuatro años}; y \textbf{1937 y 1938 enteros} ---ediciones 1669 a 1773---, el primer tramo en que los PDF traen capa de texto del proveedor: el 99\,\% de las hojas de 1937 no necesitó reconocimiento propio, y a cambio ese tramo no tiene confianza por hoja ni medición de tinta, lo que se declara como límite y no como pendiente; y \textbf{1939 y 1940 enteros}, con la misma capa y el mismo límite. Los años 1947 a 1950, barridos al cierre por pre-informe automático y sin lectura sobre la imagen, no forman parte de esas ventanas. En esas dos ventanas, y sólo ahí ---con la salvedad de los números y las páginas que faltan en cada una---, este libro puede afirmar que algo \textit{no está}. El apéndice \ref{cap:fuentes} detalla el método y enumera los ceros medidos.

\textbf{Cuarta.} Este relevamiento se equivocó más de una docena de veces, y \textbf{todas las correcciones
están incorporadas al cuerpo del texto en lugar de disimuladas}: se señala qué decía antes,
qué documento lo desmintió y qué dice ahora. Una fuente resultó pertenecer a otra provincia; una cláusula que parecía un vacío normativo resultó estar reglada en un instrumento que aún no se había leído; y un municipio que parecía carecer de reglamento de fraccionamiento resultó tenerlo sin publicar. Entre las restantes, dos merecen mención porque el libro las corrigió contra sí mismo: una fuente intermedia atribuyó a un acto administrativo cambios que su anexo original no contenía, y una observación sobre la unicidad de una omisión resultó falsa al ampliarse la serie. Se dejan asentadas porque describen el objeto: investigar sobre una normativa cuyos anexos no se publican produce, con regularidad, la apariencia de vacíos donde hay reglas, y trabajar con resúmenes en lugar de originales produce hallazgos que no existen.

El libro termina con tres piezas de naturaleza distinta, a las que se suma el capítulo prospectivo del que se habla más abajo. El capítulo~\ref{cap:conclusion} resume lo que el
relevamiento estableció y lo que no. El~\ref{cap:infraestructura} es una lectura del autor: relee lo documentado y agrega pocas piezas nuevas ---declaraciones vecinales, boletines municipales de la capital y un grupo de actos de 1935 a 1945 que sólo se desarrollan allí---,
y propone qué significa junto lo documentado. Y el~\ref{cap:plan} es, con el séptimo apéndice, lo único que este libro propone: qué habría que hacer, en qué orden y quién tendría que decidirlo. Los dos últimos van
marcados como tales, con sus condiciones de refutación, para que se los pueda discutir sin
discutir los documentos.

Los apéndices no son ilustración: son el aparato con el que se puede continuar o refutar
este trabajo. La cronología ordena los hechos desde 1690, con el grueso a partir de 1908; la matriz reúne los emprendimientos
con lo que se sabe y lo que no de cada uno; el inventario normativo transcribe lo que se pudo
obtener; el cuarto apéndice reúne, por organismo, todo lo que falta pedir; y el quinto reúne
a las personas que el archivo registra más de una vez, con la advertencia expresa de que son
coincidencias de nombre y no identificaciones. El sexto declara de qué clase de fuente sale
cada parte del libro, y qué clase de fuente no tuvo; y el octavo reconstruye, acto por acto, cómo circuló la tierra del departamento entre 1909 y 1945.

\textbf{Y hay un capítulo que mira hacia adelante, el \ref{cap:prospectiva}}, el único que proyecta ---el \ref{cap:plan} propone, pero no pronostica---: somete a prueba la palabra del título ---\textit{dispositivo}--- contra lo que los capítulos anteriores documentan, y de esa prueba deriva qué dinámicas están en curso. \textbf{Argumenta más de lo que documenta y lo dice al comenzar.} Quien quiera sólo lo verificable puede saltearlo sin perder nada de lo demás.

\textbf{Y un capítulo, el \ref{cap:ausencias}, hace lo mismo con el propio libro}: inventaría las dimensiones de la vida del departamento que este relevamiento no cubre ---la escuela, la salud, el comercio, el transporte, la iglesia, las mujeres---, reúne lo poco que el archivo dejó de cada una y dice qué haría falta. Está publicado a propósito, para que la lista de lo que falta sea pública antes de que exista el trabajo que la llene.

\textbf{Y hay un séptimo apéndice que es de otra naturaleza y conviene avisarlo acá.} Propone texto
normativo: seis ordenanzas, quince leyes y cinco actos del Ejecutivo, con su articulado. Nace
de una constatación del último capítulo --- que en los ocho puntos donde este territorio se
decide, nadie está obligado a nada --- y traduce cada uno de esos márgenes en un deber con
plazo. \textbf{Ninguna de las veintiséis propuestas manda decidir de un modo determinado}: mandan
publicar, fundar, registrar y cruzar. El apéndice aclara además que son borradores de
contenido y no de técnica legislativa, y que quien los tome deberá pasarlos por asesoría. 

Las cuarenta y una láminas tampoco son ilustración. Siete son facsímiles de cartografía, cuadros y notas oficiales que provienen de una misma pieza, el incidente de cumplimiento del amparo colectivo ---Expte.\ INC 23860/1--- que analiza el capítulo~\ref{cap:amparo}: tres mapas temáticos de la Secretaría de Recursos Hídricos de mayo de 2022, la hoja de firmas de ese informe, un cuadro de informes de impacto ambiental de canteras, la nota de la autoridad minera que lo acompaña y el Anexo I del atlas del Código de Ordenamiento Territorial. Ninguna de esas siete estaba publicada en un boletín ni en un portal público. Otras tres provienen de los documentos del Plan de Manejo: su tabla de costos y dos diapositivas de la presentación del Plan Maestro de Drenaje. Once salen de fuentes públicas de consulta difícil: el Anexo II del mismo atlas tal como lo publica el municipio, el plano de mensura 000962 y un detalle de su recuadro vial, cuatro capturas del visor cartográfico provincial, tres mapas de la red de gas del ENARGAS y el organigrama de un ministerio provincial publicado como anexo del Boletín. Las veinte restantes son de elaboración propia ---once mapas, el de ubicación del departamento, el de las fincas del archivo sobre el catastro vigente, el de la numeración de los planos de mensura, el de las líneas de ribera con coordenadas publicadas en los edictos, el de las líneas de ribera sobre el catastro, el de formación por radio censal, el del ordenamiento de bosques en el departamento, el del catastro que reclama la comunidad Kondorwaira, el de los actos de agua del departamento, el de dónde caen hoy los topónimos de 1909 a 1948 y el de las tres fracciones censales, y nueve gráficos: censo, padrón, caudales, la matriz de la toponimia en trece nóminas, la línea de tiempo de los cargos, la coparticipación municipal de 1947, las defensas contra crecientes de 1947 y 1949, la cobertura del Boletín Oficial que este libro leyó y lo que el Estado votó frente a lo que llegó--- y cada epígrafe dice con qué datos se hicieron y, en catorce de las veinte, con qué script: los scripts, las capas recortadas al departamento y las figuras están en el repositorio cartográfico \texttt{mapas\_caldera}, cuyo inventario detalla el apéndice \ref{cap:fuentes}. \textbf{Siete más tres más once más veinte son las cuarenta y una}, y el apéndice \ref{cap:fuentes} declara con qué se produjo cada una de las veinte propias, cuáles catorce son reproducibles por cualquiera ---script publicado, capas en un \textit{Release} y la cadena probada de punta a punta--- y cuáles seis no lo son todavía. \textbf{Se reproducen porque el argumento de
este libro sobre la opacidad se sostiene o cae en documentos como estos}, y el lector tiene
que poder mirarlos. El epígrafe de cada una dice de qué hoja salió; los retoques se limitaron a recortar márgenes, girar alguna hoja y ampliar un detalle, y cada epígrafe lo declara. Donde la reproducción no alcanza para medir --- y en
varias no alcanza --- se dice en el mismo lugar. \textbf{Y la imagen de cubierta no es una de las cuarenta y una}: es una composición del autor ---dibujo a mano, digitalización, \textsc{corel draw} y una pasada de optimización con un modelo generativo de imagen--- que enuncia el objeto del libro y no documenta nada. El apéndice \ref{cap:fuentes} lo declara con su cadena de producción, y deja constancia de que ninguna de las cuarenta y una láminas pasó por un modelo generativo.

Y hay algo que este relevamiento encontró tarde y que conviene saber desde el principio: \textbf{la zona alta de la cuenca es territorio de ocupación tradicional de una comunidad kolla}, relevado por el Estado nacional en 2011 y ratificado en 2023. Todo lo que este libro documenta ---los partidos de 1909, las fincas, los loteos--- se superpone a una capa anterior que el Boletín Oficial no registra.

Dos declaraciones de interés, que el cuerpo desarrolla donde corresponde y que conviene conocer desde aquí. La línea de investigación sobre la sociedad titular del loteo El Durazno llegó a este trabajo por encargo de un tercero interesado, con nombre puesto de antemano; se la siguió con los mismos estándares que al resto (capítulo~\ref{cap:loteo}). Y la consultora del sector Educación del plan estratégico provincial que el libro analiza es hermana del autor (capítulo~\ref{cap:trabajo}).\par\medskip El libro está cerrado como relevamiento documental y abierto como investigación. No contiene entrevistas propias: todo testimonio consignado proviene de prensa y se atribuye como tal, salvo dos conversaciones informales con vecinos, que el capítulo~\ref{cap:poblacion} registra como testimonio oral. Esa carencia se declara aquí, se retoma en el cierre y resulta ser --- según se argumenta allí --- la misma carencia que impide llevar la segunda tesis más lejos.
